\documentclass[border=10pt]{standalone}
\usepackage[utf8]{inputenc}
\usepackage[T1]{fontenc}
\usepackage[french]{babel}
\usepackage{tikz}
\usepackage{xcolor}
\usetikzlibrary{shapes.geometric,arrows.meta,positioning,calc,fit,backgrounds}

% ─── Couleurs ─────────────────────────────────────────────────
\definecolor{primary}{RGB}{30,58,95}
\definecolor{admincol}{RGB}{109,40,217}
\definecolor{teachcol}{RGB}{5,150,105}
\definecolor{parentcol}{RGB}{217,119,6}
\definecolor{sysback}{RGB}{245,248,255}
\definecolor{rootcol}{RGB}{220,38,38}

% ─── Styles ───────────────────────────────────────────────────
\tikzset{
  uca/.style={ellipse,draw=admincol,fill=white,line width=0.8pt,
    text width=2.4cm,align=center,font=\fontsize{6.5}{8}\selectfont,inner sep=2pt,minimum height=0.7cm},
  ucroot/.style={ellipse,draw=rootcol,fill=white,line width=0.8pt,
    text width=2.4cm,align=center,font=\fontsize{6.5}{8}\selectfont,inner sep=2pt,minimum height=0.7cm},
  uct/.style={ellipse,draw=teachcol,fill=white,line width=0.8pt,
    text width=2.4cm,align=center,font=\fontsize{6.5}{8}\selectfont,inner sep=2pt,minimum height=0.7cm},
  ucp/.style={ellipse,draw=parentcol,fill=white,line width=0.8pt,
    text width=2.4cm,align=center,font=\fontsize{6.5}{8}\selectfont,inner sep=2pt,minimum height=0.7cm},
  ucs/.style={ellipse,draw=primary!80,fill=primary!5,line width=1pt,
    text width=2.4cm,align=center,font=\fontsize{7}{9}\selectfont\bfseries,inner sep=2pt,minimum height=0.7cm},
  assoc/.style={draw=gray!50,line width=0.55pt},
  % Generalisation UML : fleche vide (triangle creux)
  gen/.style={draw=gray!70,line width=1pt,->,>=latex},
  incl/.style={draw=blue!60,dashed,line width=0.5pt,->,>=stealth},
  extn/.style={draw=orange!60,dashed,line width=0.5pt,->,>=stealth},
}

% ─── Macro acteur (bonhomme) ──────────────────────────────────
\newcommand{\actor}[4]{%  {x}{y}{label}{color}
  \begin{scope}[shift={(#1,#2)}]
    \draw[fill=#4!20,draw=#4,line width=1pt] (0,0.38) circle(0.19);
    \draw[#4,line width=1pt](0,0.19)--(0,-0.24);
    \draw[#4,line width=1pt](-0.24,0.02)--(0.24,0.02);
    \draw[#4,line width=1pt](0,-0.24)--(-0.2,-0.54);
    \draw[#4,line width=1pt](0,-0.24)--(0.2,-0.54);
    \node[below,font=\fontsize{7.5}{9}\selectfont\bfseries,text=#4,
      text width=2.2cm,align=center] at(0,-0.62){#3};
  \end{scope}
}

% ─── Macro acteur abstrait ────────────────────────────────────
\newcommand{\absactor}[4]{%  {x}{y}{label}{color}
  \begin{scope}[shift={(#1,#2)}]
    \draw[fill=#4!20,draw=#4,line width=1.2pt,dashed] (0,0.38) circle(0.19);
    \draw[#4,line width=1pt,dashed](0,0.19)--(0,-0.24);
    \draw[#4,line width=1pt,dashed](-0.24,0.02)--(0.24,0.02);
    \draw[#4,line width=1pt,dashed](0,-0.24)--(-0.18,-0.54);
    \draw[#4,line width=1pt,dashed](0,-0.24)--(0.18,-0.54);
    \node[below,font=\fontsize{7.5}{9}\selectfont\bfseries\itshape,text=#4,
      text width=2.4cm,align=center] at(0,-0.62){#3\\{\fontsize{5.5}{7}\selectfont(abstrait)}};
  \end{scope}
}

\begin{document}
\begin{tikzpicture}

%%=========================================================
%%  FRONTIERE DU SYSTEME
%%=========================================================
\draw[rounded corners=8pt,fill=sysback,draw=primary!30,line width=2pt]
  (0.3,-20.5) rectangle (25.5,2.8);
\node[anchor=north west,font=\large\bfseries,text=primary] at (0.5,2.8)
  {Portail Num\'erique --- \'Ecole Les G\'enies};

%%=========================================================
%%  UC PARTAGEE
%%=========================================================
\node[ucs] (login) at (13.5,2.0) {Se connecter};

%%=========================================================
%%  HIERARCHIE D'ACTEURS
%%=========================================================

%% Acteur abstrait : Administrateur
\absactor{-5.0}{-4.5}{Administrateur}{admincol}

%% Acteurs fils avec arbre de generalisation
%% Arbre : ligne verticale centrale + branches

%% Directeur
\actor{-5.0}{ 0.5}{Directeur}{admincol}
%% Fondateur
\actor{-5.0}{-2.0}{Fondateur}{admincol}
%% Intendant
\actor{-5.0}{-4.5}{Intendant}{admincol}  %% meme y que abstrait -- on decale
%% => on repositionne
%% On va mettre Administrateur (abstrait) au centre de ses fils

%% Repositionnement propre :
%% fils à gauche : x=-6.5
%% Administrateur abstrait : x=-3.5, entre les fils et le systeme

\end{tikzpicture}

%% ---- Repartir from scratch avec bon layout ----

\begin{tikzpicture}

%%=========================================================
%%  FRONTIERE
%%=========================================================
\draw[rounded corners=8pt,fill=sysback,draw=primary!30,line width=2pt]
  (0.2,-20.8) rectangle (25.6,2.5);
\node[anchor=north west,font=\large\bfseries,text=primary] at (0.4,2.5)
  {Portail Num\'erique --- \'Ecole Les G\'enies};

%%=========================================================
%%  UC PARTAGEE : Se connecter
%%=========================================================
\node[ucs] (login) at (3.5,1.8) {Se connecter};

%%=========================================================
%%  HIERARCHIE ADMIN (gauche)
%%
%%  Arbre :  Administrateur (abstrait) est le noeud central.
%%  Fils : Directeur, Fondateur, Intendant, Administration, Root
%%  Fleches de generalisation : fils --> Administrateur
%%=========================================================

%%  Acteur abstrait au milieu vertical de ses fils
\absactor{-4.8}{-4.5}{Administrateur}{admincol}

%%  Fils (plus a gauche)
\actor{-7.5}{ 0.0}{Directeur}{admincol}
\actor{-7.5}{-2.0}{Fondateur}{admincol}
\actor{-7.5}{-4.5}{Intendant}{admincol}
\actor{-7.5}{-7.0}{Administration}{admincol}
\actor{-7.5}{-9.0}{Root}{rootcol}

%%  Arbre generalisation : ligne verticale en x=-6.0
%%  puis branches horizontales jusqu'a chaque fils
%%  fleche vide (triangle) arrivant sur Administrateur
\draw[gray!60,line width=0.9pt] (-6.0, 0.0) -- (-6.0,-9.0); %% tronc vertical
\foreach \ay in {0.0,-2.0,-4.5,-7.0,-9.0}
  \draw[gray!60,line width=0.9pt] (-7.1,\ay) -- (-6.0,\ay); %% branche
\draw[gen] (-6.0,-4.5) -- (-5.4,-4.5); %% fleche triangle vers Administrateur

%%  Acteur Enseignant
\actor{-4.8}{-13.5}{Enseignant}{teachcol}

%%  Acteur Parent
\actor{-4.8}{-18.0}{Parent}{parentcol}

%%=========================================================
%%  ASSOCIATIONS vers Se connecter
%%=========================================================
\draw[assoc] (-4.4,-4.5+0.19) to[out=45,in=180] (login.west);
\draw[assoc] (-4.4,-13.5+0.19) to[out=30,in=210] (login.south west);
\draw[assoc] (-4.4,-18.0+0.19) to[out=25,in=230] (login.south);

%%=========================================================
%%  BANDES PAR SOUS-ACTEUR (admin)
%%=========================================================

%% ─── DIRECTEUR (bande y=-0.5 a -2.0) ──────────────────
\draw[rounded corners=4pt,fill=admincol!5,draw=admincol!25,line width=0.6pt]
  (0.5,0.5) rectangle (25.3,-1.7);
\node[anchor=west,font=\fontsize{7}{8}\selectfont\bfseries\itshape,text=admincol!80]
  at (0.7,0.2) {Directeur};

\node[uca](dir1) at ( 2.5,-0.7) {G\'erer les classes};
\node[uca](dir2) at ( 5.5,-0.7) {G\'erer les salles};
\node[uca](dir3) at ( 8.5,-0.7) {Affecter un titulaire};
\node[uca](dir4) at (11.5,-0.7) {Consulter notes \& moy.};
\node[uca](dir5) at (14.5,-0.7) {G\'en\'erer bulletin PDF};
\node[uca](dir6) at (17.5,-0.7) {G\'erer la discipline};
\node[uca](dir7) at (20.5,-0.7) {Envoyer un message};
\node[uca](dir8) at (23.5,-0.7) {Planifier une \'epreuve};

\foreach \n in {dir1,dir2,dir3,dir4,dir5,dir6,dir7,dir8}
  \draw[assoc] (-7.1, 0.0) to[out=0,in=180] (\n.west);

%% ─── FONDATEUR (bande y=-2.2 a -3.7) ──────────────────
\draw[rounded corners=4pt,fill=admincol!5,draw=admincol!25,line width=0.6pt]
  (0.5,-2.1) rectangle (25.3,-3.8);
\node[anchor=west,font=\fontsize{7}{8}\selectfont\bfseries\itshape,text=admincol!80]
  at (0.7,-2.4) {Fondateur};
\node[font=\fontsize{6}{7}\selectfont\itshape,text=gray] at (6.0,-2.4)
  {(h\'erite de toutes les fonctions du Directeur)};

\node[uca](fon1) at ( 3.5,-3.0) {Cr\'eer un utilisateur};
\node[uca](fon2) at ( 7.5,-3.0) {R\'einitialiser MDP};
\node[uca](fon3) at (11.5,-3.0) {Configurer le syst\`eme};
\node[uca](fon4) at (15.5,-3.0) {D\'efinir l'ann\'ee acad.};
\node[uca](fon5) at (19.5,-3.0) {G\'erer les trimestres};

\foreach \n in {fon1,fon2,fon3,fon4,fon5}
  \draw[assoc] (-7.1,-2.0) to[out=0,in=180] (\n.west);

%% ─── INTENDANT (bande y=-4.0 a -5.7) ──────────────────
\draw[rounded corners=4pt,fill=admincol!5,draw=admincol!25,line width=0.6pt]
  (0.5,-4.2) rectangle (25.3,-5.8);
\node[anchor=west,font=\fontsize{7}{8}\selectfont\bfseries\itshape,text=admincol!80]
  at (0.7,-4.5) {Intendant};

\node[uca](int1) at ( 3.5,-5.0) {Enregistrer paiement};
\node[uca](int2) at ( 7.5,-5.0) {Consulter les impay\'es};
\node[uca](int3) at (11.5,-5.0) {G\'en\'erer relev\'e fin.};
\node[uca](int4) at (15.5,-5.0) {Suivre scolarit\'e};

\foreach \n in {int1,int2,int3,int4}
  \draw[assoc] (-7.1,-4.5) to[out=0,in=180] (\n.west);

%% ─── ADMINISTRATION (bande y=-6.0 a -8.0) ─────────────
\draw[rounded corners=4pt,fill=admincol!5,draw=admincol!25,line width=0.6pt]
  (0.5,-6.2) rectangle (25.3,-8.2);
\node[anchor=west,font=\fontsize{7}{8}\selectfont\bfseries\itshape,text=admincol!80]
  at (0.7,-6.5) {Administration};

\node[uca](adm1) at ( 2.8,-7.2) {Inscrire un \'el\`eve};
\node[uca](adm2) at ( 6.2,-7.2) {Modifier un \'el\`eve};
\node[uca](adm3) at ( 9.6,-7.2) {D\'esactiver un \'el\`eve};
\node[uca](adm4) at (13.0,-7.2) {Saisir une absence};
\node[uca](adm5) at (16.4,-7.2) {Saisir une sanction};
\node[uca](adm6) at (19.8,-7.2) {Rechercher un \'el\`eve};
\node[uca](adm7) at (23.2,-7.2) {Consulter dossier \'el.};

\foreach \n in {adm1,adm2,adm3,adm4,adm5,adm6,adm7}
  \draw[assoc] (-7.1,-7.0) to[out=0,in=180] (\n.west);

%% ─── ROOT (bande y=-8.4 a -10.2) ──────────────────────
\draw[rounded corners=4pt,fill=rootcol!5,draw=rootcol!25,line width=0.6pt]
  (0.5,-8.4) rectangle (25.3,-10.2);
\node[anchor=west,font=\fontsize{7}{8}\selectfont\bfseries\itshape,text=rootcol!80]
  at (0.7,-8.7) {Root (Syst\`eme)};
\node[font=\fontsize{6}{7}\selectfont\itshape,text=gray] at (8.5,-8.7)
  {(acc\`es syst\`eme complet)};

\node[ucroot](rt1) at ( 3.5,-9.3) {Acc\'eder aux journaux};
\node[ucroot](rt2) at ( 7.5,-9.3) {G\'erer droits acc\`es};
\node[ucroot](rt3) at (11.5,-9.3) {Sauvegarder BDD};
\node[ucroot](rt4) at (15.5,-9.3) {Restaurer BDD};
\node[ucroot](rt5) at (19.5,-9.3) {Purger les donn\'ees};
\node[ucroot](rt6) at (23.0,-9.3) {Configurer syst\`eme};

\foreach \n in {rt1,rt2,rt3,rt4,rt5,rt6}
  \draw[assoc] (-7.1,-9.0) to[out=0,in=180] (\n.west);

%%=========================================================
%%  ENSEIGNANT (bande y=-11.2 a -15.5)
%%=========================================================
\draw[rounded corners=4pt,fill=teachcol!5,draw=teachcol!25,line width=0.6pt]
  (0.5,-11.0) rectangle (25.3,-15.5);
\node[anchor=west,font=\fontsize{7}{8}\selectfont\bfseries\itshape,text=teachcol!80]
  at (0.7,-11.3) {Enseignant};

\node[uct](te1) at ( 3.0,-12.2) {Consulter son tableau de bord};
\node[uct](te2) at ( 7.5,-12.2) {Lister ses \'el\`eves};
\node[uct](te3) at (12.0,-12.2) {Voir profil d'un \'el\`eve};
\node[uct](te4) at (16.5,-12.2) {Saisir absences du jour};
\node[uct](te5) at (21.0,-12.2) {Consulter son planning};

\node[uct](te6) at ( 5.0,-14.0) {Consulter la liste des notes};
\node[uct](te7) at (10.0,-14.0) {Consulter les sanctions};
\node[uct](te8) at (15.0,-14.0) {Suivre les pr\'esences};

\foreach \n in {te1,te2,te3,te4,te5,te6,te7,te8}
  \draw[assoc] (-4.4,-13.5+0.19) to[out=0,in=180] (\n.west);

%% include : Voir profil <<include>> Lister eleves
\draw[incl] (te3.north) to[out=100,in=80] (te2.north)
  node[midway,above,font=\fontsize{5.5}{6}\selectfont,text=blue!70]{\textit{\guillemotleft include\guillemotright}};

%%=========================================================
%%  PARENT (bande y=-15.8 a -20.5)
%%=========================================================
\draw[rounded corners=4pt,fill=parentcol!5,draw=parentcol!25,line width=0.6pt]
  (0.5,-16.0) rectangle (25.3,-20.5);
\node[anchor=west,font=\fontsize{7}{8}\selectfont\bfseries\itshape,text=parentcol!80]
  at (0.7,-16.3) {Parent};

\node[ucp](pa1) at ( 2.5,-17.2) {Consulter tableau de bord};
\node[ucp](pa2) at ( 6.5,-17.2) {Consulter notes de l'enfant};
\node[ucp](pa3) at (10.5,-17.2) {Consulter les bulletins};
\node[ucp](pa4) at (14.5,-17.2) {Imprimer bulletin PDF};
\node[ucp](pa5) at (18.5,-17.2) {Lire les messages};
\node[ucp](pa6) at (22.5,-17.2) {Consulter les paiements};

\node[ucp](pa7) at ( 4.5,-19.0) {Changer d'enfant (onglet)};
\node[ucp](pa8) at ( 9.5,-19.0) {Contacter l'administration};
\node[ucp](pa9) at (14.5,-19.0) {Acc\'eder au portail};

\foreach \n in {pa1,pa2,pa3,pa4,pa5,pa6,pa7,pa8,pa9}
  \draw[assoc] (-4.4,-18.0+0.19) to[out=0,in=180] (\n.west);

%% extend : Imprimer PDF <<extend>> Consulter bulletins
\draw[extn] (pa4.north) to[out=95,in=85] (pa3.north)
  node[midway,above,font=\fontsize{5.5}{6}\selectfont,text=orange!80]{\textit{\guillemotleft extend\guillemotright}};

%%=========================================================
%%  LEGENDE
%%=========================================================
\begin{scope}[shift={(16.5,2.0)}]
  \draw[rounded corners=3pt,fill=white,draw=gray!30,line width=0.6pt]
    (-0.3,0.35) rectangle (9.2,-1.1);
  \node[font=\fontsize{7.5}{9}\selectfont\bfseries,text=primary] at (4.5,0.15)
    {L\'egende};
  \draw[assoc]  (0,  -0.2) -- (1.0,-0.2);
  \node[right,font=\fontsize{6.5}{7.5}\selectfont] at (1.0,-0.2) {Association};
  \draw[gen]    (3.5,-0.2) -- (4.5,-0.2);
  \node[right,font=\fontsize{6.5}{7.5}\selectfont] at (4.5,-0.2) {G\'en\'eralisation (h\'eritage)};
  \draw[incl]   (0,  -0.7) -- (1.0,-0.7);
  \node[right,font=\fontsize{6.5}{7.5}\selectfont] at (1.0,-0.7) {\textit{\guillemotleft include\guillemotright}};
  \draw[extn]   (3.5,-0.7) -- (4.5,-0.7);
  \node[right,font=\fontsize{6.5}{7.5}\selectfont] at (4.5,-0.7) {\textit{\guillemotleft extend\guillemotright}};
\end{scope}

%%=========================================================
%%  TITRE BAS
%%=========================================================
\node[font=\fontsize{7}{8}\selectfont\itshape,text=gray!60,anchor=south east]
  at (25.4,-20.8)
  {\'Ecole Les G\'enies --- Portail Num\'erique v1.0 --- Diagramme UML Cas d'Utilisation --- Juillet 2026};

\end{tikzpicture}
\end{document}
